\section{Background and related work}\label{sec:background}

\subsection{Customer baselines in wholesale demand response}\label{sec:background:rules}

Wholesale programs settle on a baseline built from recent like days: the rule averages each hour over selected prior weekdays, and metered load from the hours before the event window may then adjust the result. Table~\ref{tab:methods} summarizes the methods, and Supplement Table~\ref{tabS:methodsfull} gives each rule's source, vintage and our deviation.

PJM Interconnection (PJM) takes five weekdays within 45 days, drops the lowest and averages the highest four. Its same-day adjustment is additive: the rule adds the mean gap between metered and baseline load over three hours that end one hour before the event to every baseline hour \citep{pjm_manual11,pjm_attk}. Neither the tariff text (Attachment K Appendix, s3.3A.3(a), 2015 version) nor Manual~11 (revision 137, s10.4.2) states a cap; earlier rule vintages were not checked. The ten-in-ten rule of the California Independent System Operator (CAISO) averages the ten most recent like days within 45 days and scales the profile by the ratio of actual to baseline load in the same three hours, clipped at twenty percent \citep{caiso_tariff}. The New York Independent System Operator (NYISO) and ISO New England (ISO-NE) follow the same pattern with their own selection windows and adjustment limits \citep{nyiso_manual7,isone_mr1}. Research evaluations add day matching (DM), weather matching (WM) and regression baselines with adjustments capped at 20 and 40 percent \citep{granderson2021,mathieu2011}.

Two deviations apply: each rule comes from the document current at drafting while the data cover 2016 and 2017, and we approximate the five-minute ISO-NE rule on hourly BDG2 data. Because implementation choices alone moved estimated shed by 20 to 30 kW at commercial facilities in one field study \citep{addy2015}, we specify every rule in full and freeze it.

The rules differ in the data stream they read. Unadjusted rules read only days before the event (I-hist). Adjusted rules also read the event-day hours before the window (I-morn), which improves hot-day accuracy and exposes them to load placed in those hours. Donor baselines read other buildings' same-day load (I-donor).

\subsection{Accuracy evaluations and statistical practice}\label{sec:background:accuracy}

Accuracy studies of non-residential baselines use proxy event days because real events carry no truth. Coughlin et al.\ \citep{coughlin2009} and Granderson et al.\ \citep{granderson2021} report that same-day adjustment improves accuracy on such days; the latter also report that unadjusted rules under-predict peak-window load. Our error is $e=B-L$, which is positive for over-payment. The normalized mean bias error of these studies \citep{granderson2016} has the opposite sign, so Supplement Table~\ref{tab:bridge} adds a column in their convention. Wider model families appear in whole-building models and reviews \citep{granderson2015,valentini2022,malhotra2024}.

Statistical inference is rarely reported. The cited studies give medians, box plots or mean errors, usually without intervals or tests \citep{granderson2021,hatton2016,lee2019}; one compares methods with t-tests across customers \citep{wang2018spm}. Commercial data are proprietary or held by the authors \citep{granderson2021,mathieu2011cdc}, one data set per study, mostly at 15-minute resolution \citep{coughlin2009,mathieu2011}.

\subsection{Comparison groups and synthetic control}\label{sec:background:donors}

A comparison group, called a control group in the literature, estimates the counterfactual from non-participants observed on the event day. Selected or randomized control groups were more accurate than day and weather matching in residential trials \citep{hatton2016,wang2018spm,lee2019}, and California's baseline working group proposed them as a settlement option \citep{caiso_bawg2017}. The synthetic control method (SCM) weights donors to match the target before the event \citep{abadie2010,abadie2021}. Ridge forms drop the sign and intercept restrictions \citep{benmichael2021,doudchenko2016}, and a 2026 preprint applies one to residential data \citep{sievers2026}. These studies are residential or simulated, test superiority and not equivalence, and do not stress the donor baseline with injected shaping. Our CG-adj is PJM-adj with a donor ratio, not an independent donor method. SCM-ridge reads no target hour in the shaped set.

\subsection{Pre-event shaping, gaming and repeated dispatch}\label{sec:background:shaping}

Administrative baselines invite gaming \citep{chao2011}. In one simulation, a rational customer under-consumes on DR days and over-consumes on other days to raise a high-X-of-Y baseline \citep{wang2022manipulation}. Gabaldon et al.\ separate licit pre-conditioning from gaming with a physically based HVAC model \citep{gabaldon2021}, and Afentoulis and Vagropoulos find that two of four baseline designs let a simulated electric-vehicle aggregator raise revenue \citep{afentoulis2025}. Field and policy documents report qualitatively that pre-cooling or rebound in the adjustment hours distorts a same-day adjusted baseline \citep{coughlin2009,braithwait2010,recurve2022caiso}, and the California working group proposed a buffer before the adjustment window \citep{caiso_bawg2017}. Todd et al.\ show that spillover of event-day reductions onto non-event days biases baselines, with a sign opposite to that of pre-event inflation \citep{todd2019}. We use the term pre-event load increase and assert no behavior, because the rules cannot separate licit pre-conditioning from gaming.

Baseline erosion has a signed meaning in the literature: the estimated curtailment drifts as event days enter the lookback \citep{enernoc2009naesb,caiso_bawg2017}. A 2011 analysis for PJM left the effect of frequent settlements unanalyzed \citep{kema2011pjm}, and a California evaluation found that nine consecutive event days pushed residential baseline days back by nearly two weeks \citep{dsa2023elrp}. Granderson et al.\ slid the lookback on non-DR data and report growing dispersion \citep{granderson2023}, and Wang et al.\ attribute the difficulty of high-frequency DR to scarce reference days in simulated homes \citep{wang2026physics}. We use erosion only in the signed sense. The quantity tested here, the growth of absolute estimation error with dispatch count, is consistent in sign with that dispersion growth.

\subsection{What remains open}\label{sec:background:open}

Three questions remain. (i) How large an injected four-hour pre-event increase does PJM-adj tolerate before a baseline that does not read those hours becomes more accurate? (ii) How does absolute estimation error grow with dispatch count for a rule and for a synthetic control, site by site, with exact inference? (iii) Does a comparison group add anything to PJM-adj at the site level? The computing-side work reviewed in Section~\ref{sec:intro} does not estimate the settlement counterfactual and addresses none of them.
